\section*{الفصل \ref{ch:g4:breathing} --- التنفس}

\begin{solution}{exo:g4:breathing:1}
\omterm{def:g4:breathing:movements}{الشهيق}: يتسع الصدر فيدخل الهواء. \omterm{def:g4:breathing:movements}{والزفير}: يرتخي الصدر فيخرج
الهواء.
\end{solution}

\begin{solution}{exo:g4:breathing:2}
الأنف (أو الفم)، \omterm{def:g4:breathing:organs}{فالقصبة الهوائية}، فالفرعان، فالرئتان، وأخيرًا
ملايين الجيوب الهوائية الصغيرة الملفوفة بالأوعية الدموية.
\end{solution}

\begin{solution}{exo:g4:breathing:3}
نحو $15$ مرة في الدقيقة. وفي العدو يرتفع إلى $40$ أو أكثر، وأعمق
أيضًا --- \omterm{def:g3:bones-and-muscles:muscle}{فالعضلات} الكادحة تحتاج إلى الأكسجين أسرع، فتحمّله الرئتان
أسرع.
\end{solution}

\begin{solution}{exo:g4:breathing:4}
يأخذ الأكسجين ويردّ ثاني أكسيد الكربون. والتبادل يحدث في جيوب
الرئتين الهوائية الصغيرة، عبر جدرانها الرقيقة، مع \omterm{prop:g4:journey-of-food:absorb}{الدم}.
\end{solution}

\begin{solution}{exo:g4:breathing:5}
هواء \omterm{def:g4:breathing:movements}{الزفير} أدفأ وأرطب (وهو ضباب المرآة)، وأفقر أكسجينًا وأغنى
ثاني أكسيد كربون.
\end{solution}

\begin{solution}{exo:g4:breathing:6}
لأن تبادل الهواء \omterm{prop:g4:journey-of-food:absorb}{والدم} يحتاج إلى مساحة: فجدران ملايين الجيوب
تتجمع في نصف ملعب تنس من مساحة اللقاء مطوية في الصدر. \omterm{def:g4:journey-of-food:tract}{والمعي
الدقيق} يلعب الحيلة نفسها بطيّاته لامتصاص \omterm{def:g4:journey-of-food:digestion}{المغذيات}.
\end{solution}

\begin{solution}{exo:g4:breathing:7}
الأنف يدفّئ الهواء ويرطّبه ويصفّيه في طريقه إلى الداخل --- فيبلغ
الهواءُ الباردُ الجافُّ المغبَّرُ الرئتين الرقيقتين مهيَّأً.
\end{solution}

\begin{solution}{exo:g4:breathing:8}
الجواب مفتوح. وفي العادة نحو $15$ في الراحة و$25$--$40$ بعد
القفز: \omterm{def:g3:bones-and-muscles:muscle}{فالعضلات} العاملة طلبت أكسجينًا أكثر، فأسرع \omterm{def:g4:breathing:movements}{التنفس} وعمق
ليوفّره.
\end{solution}

\begin{solution}{exo:g4:breathing:9}
على أن الجسم يخزن الغذاء بسخاء (ما يكفي ساعات أو أيامًا) لكنه لا
يحتفظ بمخزون من الأكسجين تقريبًا --- بالكاد ما يكفي دقيقة. فوجب
توصيل الأكسجين باستمرار، نفَسًا بنفَس، ولهذا لا يستطيع \omterm{def:g4:breathing:movements}{التنفس} أن
يتوقف طويلًا.
\end{solution}

\begin{solution}{exo:g4:breathing:10}
خمسة وعشرون متنفسًا ظلوا يأخذون الأكسجين من هواء الغرفة ويردّون
ثاني أكسيد الكربون طول الصباح: فصار هواؤها أفقر بالأول وأغنى
بالثاني. والوصفة هي القاعدة الأولى: افتح النوافذ وهوّ الغرفة.
\end{solution}

\begin{solution}{exo:g4:breathing:11}
كلاهما موضع لقاء \omterm{prop:g4:journey-of-food:absorb}{الدم} والأكسجين: سطوح هائلة رقيقة مبطنة \omterm{prop:g4:journey-of-food:absorb}{بالدم}
يعبر عندها الأكسجين إلى الداخل وثاني أكسيد الكربون إلى الخارج ---
فالغلاصم تأخذ الأكسجين المذاب في الماء، والرئتان تأخذانه من
الهواء. وفي الهواء تنهار سطوح غلاصم السلمون الريشية ويلتصق بعضها
ببعض: فتنكمش مساحة العبور الكبرى حتى تكاد تنعدم، ولا مساحة عبور
يعني لا أكسجين، مهما يكن الهواء حوله غنيًّا.
\end{solution}
